\section{PAPER 025: Triang\-ulationEngine — Four-Perspective Candidate Validation for Knowledge Graph Discovery}

\textbf{Author}: Bryan Alexander Buchorn  
\textbf{Affiliation}: Independent Researcher, Las Vegas, NV, USA  
\textbf{Status}: v2.51.0 (Phase 167 (STRB) COMPLETE)
\textbf{Date}: May 2, 2026

\textbf{CEREBRUM Phase 72}

---

\subsection{Abstract}

We present \textbf{Triang\-ulationEngine}, a four-perspective validation framework for \texttt{ResearchCandidate} objects in CEREBRUM's knowledge graph discovery pipeline. Inspired by trian\-gulation in navigation and qualitative research methodology [denzin\-1978\-research], the engine validates each candidate edge from four independent persp\-ectives: (P1) \textbf{reverse traversal confidence} — does the graph support the inverse relation?; (P2) \textbf{multi-strategy agreement} — do different reasoning confi\-gurations agree?; (P3) \textbf{path indep\-endence} — are the supporting paths struc\-turally independent?; (P4) \textbf{semantic type consistency} — is the relation type compatible with the entity class profile? The four perspective scores extend the \texttt{AutoApprover} feature vector from 12 to 16 dimensions, providing richer signal for the downstream logistic classifier. A diagnostic \texttt{is\_Synaptic Bridge\_candidate} flag identifies cross-community bridge proposals warranting special handling.

---

\subsection{1. Motivation: Single-Path Validation is Insuff\-icient}

Prior to Phase 72, \texttt{ResearchAgent} validated candidates by running \texttt{HypothesisEngine} once in the forward direction (A$\to$B) and checking whether supporting paths exceeded a confidence threshold. This single-perspective approach has three failure modes:

\begin{enumerate}
\item \textbf{Direct\-ionality bias}: A$\to$B may traverse well even when B$\to$A has no support, producing spurious asymmetric edges.
\item \textbf{Reasoning monoculture}: A single reasoning confi\-guration may over-weight community structure ($\beta$) and consi\-stently produce high-confidence paths for topol\-ogically close but seman\-tically unrelated entities.
\item \textbf{Dependent paths}: Multiple "independent" paths through the same bottleneck node provide weaker evidence than truly parallel routes.

\end{enumerate}
Triang\-ulation addresses all three by requiring convergent evidence across struc\-turally independent persp\-ectives.

---

\subsection{2. The Four Perspe\-ctives}

\subsubsection{P1: Reverse Traversal Confidence}

Run \texttt{HypothesisEngine} from target B to source A (inverse direction):

\begin{lstlisting}
reverse_result = hypothesis_engine.evaluate(target, source)
p1 = reverse_result.confidence
\end{lstlisting}

A genuine causal or associative relat\-ionship should exhibit non-trivial reverse traversal confidence. Spurious forward-only paths yield near-zero reverse confidence.

\subsubsection{P2: Multi-Strategy Agreement}

Run the candidate through three different \texttt{BeamTraversal} confi\-gurations (varying \texttt{beam\_width}, \texttt{probabilistic}, \texttt{max\_hop}) and compute the agreement fraction:

\begin{lstlisting}
scores = [config_A.score(A, B), config_B.score(A, B), config_C.score(A, B)]
p2 = len([s for s in scores if s > threshold]) / len(scores)
\end{lstlisting}

High agreement across confi\-gurations indicates a robust signal, not a confi\-guration-specific artifact.

\subsubsection{P3: Mean Path Indepe\-ndence}

Given the primary supporting paths \texttt{{P\_1, ..., P\_k}}, compute pairwise Jaccard distance between path node sets and average:

\begin{lstlisting}
independence_scores = [
    1 - jaccard(set(P_i.nodes), set(P_j.nodes))
    for i, j in combinations(range(len(paths)), 2)
]
p3 = mean(independence_scores) if independence_scores else 0.5
\end{lstlisting}

High indep\-endence (p3 $\to$ 1.0) means paths traverse different graph regions — stronger evidence. Low indep\-endence (p3 $\to$ 0.0) means all paths share the same bottleneck node — single point of failure.

\subsubsection{P4: Semantic Type Score}

Check relation-type and entity-class consistency using a type profile derived from existing graph edges:

\begin{lstlisting}
p4 = type_consistency(source_entity_class, target_entity_class, relation_type)
\end{lstlisting}

\begin{itemize}
\item Known-compatible type combination $\to$ p4 = 1.0
\item Known-incom\-patible $\to$ p4 = 0.0
\item Novel / unseen relation type $\to$ p4 = 0.5 (neutral, no penalty for discovery)

\end{itemize}
\subsubsection{Synaptic Bridge Candidate Flag}

\texttt{is\_Synaptic Bridge\_candidate = True} when source and target belong to communities with large structural distance (\textgreater  2 hops) and P1 $\times$ P2 $\times$ P3 product \textgreater  0.3. Synaptic Bridge candidates are high-value cross-community bridge proposals.

---

\subsection{3. Integration}

\begin{lstlisting}
from core.triangulation_engine import TriangulationEngine

engine = TriangulationEngine(hypothesis_engine, traversal, adapter)
report = engine.validate(candidate)

# report.reverse_confidence  → P1
# report.strategy_agreement  → P2
# report.mean_path_independence → P3
# report.semantic_type_score → P4
# report.is_Synaptic Bridge_candidate → bool

finding.metadata["triangulation"] = report
\end{lstlisting}

The four scores are autom\-atically incor\-porated into \texttt{AutoApprover}'s 16-feature vector when \texttt{triangulation} metadata is present.

---

\subsection{4. References}

\begin{itemize}
\item [denzin\-1978\-research] Denzin, N.K. (1978). \textit{The research act: A theoretical intro\-duction to socio\-logical methods.} McGraw-Hill.

\end{itemize}
---
\textbf{Copyright © 2026 Bryan Alexander Buchorn. All Rights Reserved.}

---
\subsection{Acknow\-ledgments}

The author gratefully ackno\-wledges the use of Claude (Anthropic) as a research assistant throughout this work. Claude assisted with mathe\-matical forma\-lization, code generation, manuscript preparation, and technical writing. All conceptual contr\-ibutions, archi\-tectural decisions, exper\-imental design, and intel\-lectual claims are solely the author's.

\textbf{Reviewed on}: May 2, 2026 for version v2.51.0
